% questions.tex —— 2014 年·天津卷·文科 · 题目内容片段（无答案）
% 由 tools/generate_projects.py 依来源总账自 2014/2014_数学真题_天津文_试题.md 生成，勿手改（改源 MD 后重跑）


\begin{gknotice}
\textbf{来源说明}：本题面依据固定源 PDF 的两页 A3 题面逐题转录，原卷印刷页码为 925—926。第 925 面为选择题 1—8、填空题 9—14，以及解答题 15、16 的题干；第 926 面为解答题 17—20。

\textbf{图形说明}：第 7 题圆内接三角形与切线、第 10 题三视图、第 11 题程序框图、第 17 题四棱锥，均由项目内 TikZ 矢量源重绘；其余各题原页无图。

\textbf{解析归属}：固定源不附答案或解析；配套答案与解析均由本库按本题面独立复算。
\end{gknotice}

\gkpart{一、选择题}

\begin{question}
\gkti{1}$i$ 是虚数单位，复数 $\dfrac{7+\mathrm i}{3+4\mathrm i}=$\mbox{（　　）}

（A）$1-\mathrm i$；（B）$-1+\mathrm i$；（C）$\dfrac{17}{25}+\dfrac{31}{25}\mathrm i$；（D）$-\dfrac{17}{7}+\dfrac{25}{7}\mathrm i$。

\end{question}
\begin{question}
\gkti{2}设变量 $x$，$y$ 满足约束条件
\nopagebreak[4]
\[\begin{cases}
x+y-2\geqslant 0,\\
x-y-2\leqslant 0,\\
y\geqslant 1,
\end{cases}\]
则目标函数 $z=x+2y$ 的最小值为\mbox{（　　）}

（A）$2$；（B）$3$；（C）$4$；（D）$5$。

\end{question}
\begin{question}
\gkti{3}已知命题 $p$：$\forall x>0$，总有 $(x+1)\mathrm e^x>1$，则 $\neg p$ 为\mbox{（　　）}

（A）$\exists x_0\leqslant 0$，使得 $(x_0+1)\mathrm e^{x_0}\leqslant 1$；（B）$\exists x_0>0$，使得 $(x_0+1)\mathrm e^{x_0}\leqslant 1$；（C）$\forall x>0$，总有 $(x+1)\mathrm e^x\leqslant 1$；（D）$\forall x\leqslant 0$，总有 $(x+1)\mathrm e^x\leqslant 1$。

\end{question}
\begin{question}
\gkti{4}设 $a=\log_2\pi$，$b=\log_{\frac12}\pi$，$c=\pi^{-2}$，则\mbox{（　　）}

（A）$a>b>c$；（B）$b>a>c$；（C）$a>c>b$；（D）$c>b>a$。

\end{question}
\begin{question}
\gkti{5}设 $\{a_n\}$ 是首项为 $a_1$，公差为 $-1$ 的等差数列，$S_n$ 为其前 $n$ 项和，若 $S_1$，$S_2$，$S_4$ 成等比数列，则 $a_1=$\mbox{（　　）}

（A）$2$；（B）$-2$；（C）$\dfrac12$；（D）$-\dfrac12$。

\end{question}
\begin{question}
\gkti{6}已知双曲线 $\dfrac{x^2}{a^2}-\dfrac{y^2}{b^2}=1$（$a>0$，$b>0$）的一条渐近线平行于直线 $l:y=2x+10$，双曲线的一个焦点在直线 $l$ 上，则双曲线的方程为\mbox{（　　）}

（A）$\dfrac{x^2}{5}-\dfrac{y^2}{20}=1$；（B）$\dfrac{x^2}{20}-\dfrac{y^2}{5}=1$；（C）$\dfrac{3x^2}{25}-\dfrac{3y^2}{100}=1$；（D）$\dfrac{3x^2}{100}-\dfrac{3y^2}{25}=1$。

\Needspace{0.4\textheight}
\end{question}
\begin{question}
\gkti{7}如图，$\triangle ABC$ 是圆的内接三角形，$\angle BAC$ 的平分线交圆于点 $D$，交 $BC$ 于 $E$，过点 $B$ 的圆的切线与 $AD$ 的延长线交于点 $F$。在上述条件下，给出下列四个结论：① $BD$ 平分 $\angle CBF$；② $FB^2=FD\cdot FA$；③ $AE\cdot CE=BE\cdot DE$；④ $AF\cdot BD=AB\cdot BF$。则所有正确结论的序号是\mbox{（　　）}
\nopagebreak[4]
\begin{center}% ============================================================
% f2014_tianjin_arts_q07_circle.tex
% 2014 天津卷（文科）第 7 题：圆内接三角形与切线
% 来源 PDF：普通高考/2014/2014天津文.pdf
% 原页：第 1 页（A3），印面 925，第 7 题题干下方
% D 为 ∠BAC 平分线与圆的交点，E=AD∩BC，F 为过 B 的切线与 AD 延长线的交点。
% 取 ∠A=80°，B、C 的弧上位置 208°、338°，此配置下结论 ①②④ 成立、③ 不成立。
% ============================================================
\begin{tikzpicture}[>=Stealth, line width=0.7pt, font=\small,
  line cap=round, line join=round, scale=1.2]
  % 圆（半径 1）
  \draw[line width=0.7pt] (0,0) circle (1);
  \coordinate (A) at (0.174,0.985);
  \coordinate (B) at (-0.883,-0.469);
  \coordinate (C) at (0.927,-0.375);
  \coordinate (D) at (0.052,-0.999);
  \coordinate (E) at (0.088,-0.419);
  \coordinate (F) at (-0.015,-2.102);

  % 三角形与内接弦
  \draw (A) -- (B) -- (C) -- cycle;
  \draw (B) -- (D) -- (C);
  % 角平分线 AD 及延长线到 F
  \draw (A) -- (D);
  \draw (D) -- (F);
  % 过 B 的切线
  \draw (B) -- (F);

  \fill (A) circle (0.9pt) node[above right=-1pt] {$A$};
  \fill (B) circle (0.9pt) node[below left=-1pt] {$B$};
  \fill (C) circle (0.9pt) node[right=-1pt] {$C$};
  \fill (D) circle (0.9pt) node[below right=-1pt] {$D$};
  \fill (E) circle (0.9pt) node[above left=-1pt] {$E$};
  \fill (F) circle (0.9pt) node[below=-1pt] {$F$};
\end{tikzpicture}
\end{center}
\nopagebreak[4]
（A）①②；（B）③④；（C）①②③；（D）①②④。

\end{question}
\begin{question}
\gkti{8}已知函数 $f(x)=\sqrt3\sin\omega x+\cos\omega x$（$\omega>0$），$x\in\mathbf R$。在曲线 $y=f(x)$ 与直线 $y=1$ 的交点中，若相邻交点距离的最小值为 $\dfrac\pi3$，则 $f(x)$ 的最小正周期为\mbox{（　　）}

（A）$\dfrac\pi2$；（B）$\dfrac{2\pi}{3}$；（C）$\pi$；（D）$2\pi$。

\end{question}
\gkpart{二、填空题}

\begin{question}
\gkti{9}某大学为了解在校本科生对参加某项社会实践活动的意向，拟采用分层抽样的方法，从该校四个年级的本科生中抽取一个容量为 $300$ 的样本进行调查。已知该校一年级、二年级、三年级、四年级的本科生人数之比为 $4:5:5:6$，则应从一年级本科生中抽取\gkblank{}名学生。

\Needspace{0.4\textheight}
\end{question}
\begin{question}
\gkti{10}已知一个几何体的三视图如图所示（单位：m），则该几何体的体积为\gkblank{}$\mathrm m^3$。
\nopagebreak[4]
\begin{center}% ============================================================
% f2014_tianjin_arts_q10_views.tex
% 2014 天津卷（文科）第 10 题：三视图
% 来源 PDF：普通高考/2014/2014天津文.pdf
% 原页：第 1 页（A3），印面 925，第 10 题题干下方
% 正视图、侧视图相同：宽 2、高 4 的矩形上加底宽 4、高 2 的三角形；
% 俯视图为半径 2 的圆，内含半径 1 的虚线圆（下部圆柱被遮挡的轮廓）。
% ============================================================
\begin{tikzpicture}[>=Stealth, line width=0.6pt, font=\small, scale=0.55,
  line cap=round, line join=round]
  % ---------- 正视图 ----------
  % 下部矩形：x∈[1,3]，高 4；上部三角形：底 x∈[0,4] 于 y=4，顶点 (2,6)
  \draw (1,0) -- (3,0) -- (3,4) -- (4,4) -- (2,6) -- (0,4) -- (1,4) -- (1,0);
  % 底宽 2（矩形）
  \draw[<->] (1,-0.5) -- (3,-0.5);
  \node[below] at (2,-0.5) {$2$};
  % 底宽 4（三角形底边）
  \draw[<->] (-0.6,-1.3) -- (4.6,-1.3);
  \draw (-0.6,-0.9) -- (-0.6,-1.7);
  \draw (4.6,-0.9) -- (4.6,-1.7);
  \node[below] at (2,-1.3) {$4$};
  \node[below] at (2,-2.2) {正视图};
  % 柱高 4、锥高 2（置于右侧）
  \draw[<->] (5.4,0) -- (5.4,4);
  \draw (5.0,0) -- (5.8,0);
  \draw (5.0,4) -- (5.8,4);
  \node[right] at (5.4,2) {$4$};
  \draw[<->] (5.4,4) -- (5.4,6);
  \draw (5.0,6) -- (5.8,6);
  \node[right] at (5.4,5) {$2$};

  % ---------- 侧视图 ----------
  \begin{scope}[shift={(9.2,0)}]
    \draw (1,0) -- (3,0) -- (3,4) -- (4,4) -- (2,6) -- (0,4) -- (1,4) -- (1,0);
    \draw[<->] (1,-0.5) -- (3,-0.5);
    \node[below] at (2,-0.5) {$2$};
    \draw[<->] (-0.6,-1.3) -- (4.6,-1.3);
    \draw (-0.6,-0.9) -- (-0.6,-1.7);
    \draw (4.6,-0.9) -- (4.6,-1.7);
    \node[below] at (2,-1.3) {$4$};
    \node[below] at (2,-2.2) {侧视图};
  \end{scope}

  % ---------- 俯视图 ----------
  \begin{scope}[shift={(2,-6.6)}]
    \draw (0,0) circle (2);
    \draw[dashed] (0,0) circle (1);
    \node[below] at (0,-2.7) {俯视图};
  \end{scope}
\end{tikzpicture}
\end{center}
\nopagebreak[4]
\Needspace{0.4\textheight}
\end{question}
\begin{question}
\gkti{11}阅读如图的框图，运行相应的程序，输出 $S$ 的值为\gkblank{}。
\nopagebreak[4]
\begin{center}% ============================================================
% f2014_tianjin_arts_q11_flowchart.tex
% 2014 天津卷（文科）第 11 题：程序框图
% 来源 PDF：普通高考/2014/2014天津文.pdf
% 原页：第 1 页（A3），印面 925，第 11 题题干下方
% S=0,n=3；循环体 S=S+(-2)^n、n=n-1；n⩽1 成立则输出 S，否则返回循环体。
% ============================================================
\begin{tikzpicture}[>=Stealth, line width=0.6pt, font=\small, scale=0.85,
  line cap=round, line join=round,
  term/.style={draw, rectangle, rounded corners=2.2pt, minimum width=1.9cm,
    minimum height=0.72cm, inner sep=2pt},
  proc/.style={draw, rectangle, minimum width=3.9cm, minimum height=0.72cm,
    inner sep=2pt},
  io/.style={draw, trapezium, trapezium left angle=68, trapezium right angle=112,
    minimum width=3.4cm, minimum height=0.78cm, inner sep=2pt},
  dec/.style={draw, diamond, aspect=2.1, minimum width=4.2cm, minimum height=1.35cm,
    inner sep=1pt}]
  \node[term] (start) at (0,0) {开始};
  \node[proc] (init) at (0,-1.45) {$S=0$，$n=3$};
  \node[proc] (body) at (0,-2.95) {$S=S+(-2)^n$};
  \node[proc] (dec2) at (0,-4.45) {$n=n-1$};
  \node[dec]  (test) at (0,-6.25) {$n\leqslant 1$？};
  \node[io]   (out)  at (0,-7.95) {输出 $S$};
  \node[term] (end)  at (0,-9.30) {结束};

  \draw[->] (start) -- (init);
  \draw[->] (init) -- (body);
  \draw[->] (body) -- (dec2);
  \draw[->] (dec2) -- (test);
  \draw[->] (test) -- (out);
  \draw[->] (out) -- (end);

  % “是”标注：判断框下沿与输出框之间，箭头右侧
  \node[right, inner sep=2pt] at (0.30,-7.15) {是};

  % “否”：从判断框左顶点引出，左侧上升后折向循环体顶边
  \draw[->] (test.west) -- (-3.05,-6.25) -- (-3.05,-2.30) -- (-1.15,-2.30)
    -- (-1.15,-2.59);
  \node[above, inner sep=2pt] at (-2.55,-6.25) {否};
\end{tikzpicture}
\end{center}
\nopagebreak[4]
\end{question}
\begin{question}
\gkti{12}函数 $f(x)=\lg x^2$ 的单调递减区间是\gkblank{}。

\end{question}
\begin{question}
\gkti{13}已知菱形 $ABCD$ 的边长为 $2$，$\angle BAD=120^\circ$，点 $E$，$F$ 分别在边 $BC$，$DC$ 上，$\overrightarrow{BC}=3\overrightarrow{BE}$，$\overrightarrow{DC}=\lambda\overrightarrow{DF}$。若 $\overrightarrow{AE}\cdot\overrightarrow{AF}=1$，则 $\lambda$ 的值为\gkblank{}。

\end{question}
\begin{question}
\gkti{14}已知函数
\nopagebreak[4]
\[f(x)=\begin{cases}
|x^2+5x+4|, & x\leqslant 0,\\
2|x-2|, & x>0,
\end{cases}\]
若函数 $y=f(x)-a|x|$ 恰有 $4$ 个零点，则实数 $a$ 的取值范围为\gkblank{}。

\end{question}
\gkpart{三、解答题}

\begin{question}
\gkti{15}某校夏令营有 $3$ 名男同学 $A$，$B$，$C$ 和 $3$ 名女同学 $X$，$Y$，$Z$，其年级情况如下表：

\begin{center}
\begin{tabular}{llll}
\toprule
 & 一年级 & 二年级 & 三年级 \\
\midrule
男同学 & $A$ & $B$ & $C$ \\
女同学 & $X$ & $Y$ & $Z$ \\
\bottomrule
\end{tabular}
\end{center}

现从这 $6$ 名同学中随机选出 $2$ 人参加知识竞赛（每人被选到的可能性相同）。

（1） 用表中字母列举出所有可能的结果；

（2） 设 $M$ 为事件“选出的 $2$ 人来自不同年级且恰有 $1$ 名男同学和 $1$ 名女同学”，求事件 $M$ 发生的概率。

\end{question}
\begin{question}
\gkti{16}在 $\triangle ABC$ 中，内角 $A$，$B$，$C$ 所对的边分别为 $a$，$b$，$c$，已知 $a-c=\dfrac{\sqrt6}{6}b$，$\sin B=\sqrt6\sin C$。

（1） 求 $\cos A$ 的值；

（2） 求 $\cos\left(2A-\dfrac\pi6\right)$ 的值。

\Needspace{0.4\textheight}
\end{question}
\begin{question}
\gkti{17}如图，四棱锥 $P\text{-}ABCD$ 的底面 $ABCD$ 是平行四边形，$BA=BD=\sqrt2$，$AD=2$，$PA=PD=\sqrt5$，$E$，$F$ 分别是棱 $AD$，$PC$ 的中点。
\nopagebreak[4]
\begin{center}% ============================================================
% f2014_tianjin_arts_q17_pyramid.tex
% 2014 天津卷（文科）第 17 题：四棱锥 P-ABCD
% 来源 PDF：普通高考/2014/2014天津文.pdf
% 原页：第 2 页（A3），印面 926，第 17 题题干下方
% 底面 ABCD 为平行四边形，E 为 AD 中点，F 为 PC 中点。
% 由 BA=BD=√2、AD=2 得 AB⊥BD，PB⊥底面 ABCD，
% 故 PB、BC、BD、EF 为不可见线（虚线）。立体示意图为平行投影，不按比例。
% ============================================================
\begin{tikzpicture}[>=Stealth, line width=0.65pt, font=\small, scale=1.15,
  line cap=round, line join=round]
  \coordinate (A) at (3.20,0);
  \coordinate (D) at (1.60,0);
  \coordinate (B) at (1.05,0.62);
  \coordinate (C) at (-0.55,0.62);
  \coordinate (P) at (1.15,2.50);
  \coordinate (E) at (2.40,0);
  \coordinate (F) at (0.30,1.56);

  % 可见棱
  \draw (P) -- (A);
  \draw (P) -- (D);
  \draw (P) -- (C);
  \draw (A) -- (B);
  \draw (A) -- (D);
  \draw (D) -- (C);

  % 不可见线
  \draw[dashed] (P) -- (B);
  \draw[dashed] (B) -- (C);
  \draw[dashed] (B) -- (D);
  \draw[dashed] (E) -- (F);

  \fill (A) circle (0.85pt) node[below right=-1pt] {$A$};
  \fill (B) circle (0.85pt) node[above left=-2pt] {$B$};
  \fill (C) circle (0.85pt) node[left=-1pt] {$C$};
  \fill (D) circle (0.85pt) node[below left=-1pt] {$D$};
  \fill (P) circle (0.85pt) node[above=-1pt] {$P$};
  \fill (E) circle (0.85pt) node[below=-1pt] {$E$};
  \fill (F) circle (0.85pt) node[left=-1pt] {$F$};
\end{tikzpicture}
\end{center}
\nopagebreak[4]
（1） 证明：$EF\parallel$ 平面 $PAB$；

（2） 若二面角 $P\text{-}AD\text{-}B$ 为 $60^\circ$，

① 证明：平面 $PBC\perp$ 平面 $ABCD$；

② 求直线 $EF$ 与平面 $PBC$ 所成角的正弦值。

\end{question}
\begin{question}
\gkti{18}设椭圆 $\dfrac{x^2}{a^2}+\dfrac{y^2}{b^2}=1$（$a>b>0$）的左、右焦点分别为 $F_1$，$F_2$，右顶点为 $A$，上顶点为 $B$，已知 $|AB|=\dfrac{\sqrt3}{2}|F_1F_2|$。

（1） 求椭圆的离心率；

（2） 设 $P$ 为椭圆上异于其顶点的一点，以线段 $PB$ 为直径的圆经过点 $F_1$，经过点 $F_2$ 的直线 $l$ 与该圆相切于点 $M$，$|MF_2|=2\sqrt2$。求椭圆的方程。

\end{question}
\begin{question}
\gkti{19}已知函数 $f(x)=x^2-\dfrac23ax^3$（$a>0$），$x\in\mathbf R$。

（1） 求 $f(x)$ 的单调区间和极值；

（2） 若对于任意的 $x_1\in(2,+\infty)$，都存在 $x_2\in(1,+\infty)$，使得 $f(x_1)\cdot f(x_2)=1$，求 $a$ 的取值范围。

\end{question}
\begin{question}
\gkti{20}已知 $q$ 和 $n$ 均为给定的大于 $1$ 的自然数。设集合 $M=\{0,1,2,\cdots,q-1\}$，集合 $A=\{x\mid x=x_1+x_2q+\cdots+x_nq^{n-1},\ x_i\in M,\ i=1,2,\cdots,n\}$。

（1） 当 $q=2$，$n=3$ 时，用列举法表示集合 $A$；

（2） 设 $s$，$t\in A$，$s=a_1+a_2q+\cdots+a_nq^{n-1}$，$t=b_1+b_2q+\cdots+b_nq^{n-1}$，其中 $a_i$，$b_i\in M$，$i=1,2,\cdots,n$。证明：若 $a_n<b_n$，则 $s<t$。

\end{question}